% =====================================================================
%  CV - Menra Wedwang Romial (published as Romial Menra)
%  Build:  cd cv && latexmk -pdf resume_romial.tex
%  Then copy resume_romial.pdf to ../web/static/pdf/resume_romial.pdf
% =====================================================================
\documentclass[10pt,a4paper]{article}

\usepackage[T1]{fontenc}
\usepackage[utf8]{inputenc}
\usepackage{lmodern}
\usepackage[scaled=0.95]{helvet}
\renewcommand{\familydefault}{\sfdefault}
\usepackage[a4paper,margin=1.6cm,top=1.3cm,bottom=1.4cm]{geometry}
\usepackage{xcolor}
\usepackage{titlesec}
\usepackage{enumitem}
\usepackage{tabularx}
\usepackage{fontawesome5}
\usepackage[hidelinks]{hyperref}

% ---- Colours ---------------------------------------------------------
\definecolor{accent}{HTML}{C8402A}
\definecolor{muted}{HTML}{5A5048}

\hypersetup{
  pdftitle={CV - Menra Wedwang Romial},
  pdfauthor={Menra Wedwang Romial (Romial Menra)},
  pdfkeywords={Menra Romial, Romial Menra, cloud computing, Kubernetes, energy-aware computing, Intel RAPL, power capping}
}

% ---- Layout helpers --------------------------------------------------
\setlength{\parindent}{0pt}
\pagestyle{empty}

\titleformat{\section}{\large\bfseries\color{accent}}{}{0pt}{\MakeUppercase}[\vspace{-0.4ex}\rule{\linewidth}{0.6pt}]
\titlespacing*{\section}{0pt}{1.6ex}{1.1ex}

\setlist[itemize]{leftmargin=1.1em, itemsep=0.15ex, topsep=0.3ex, parsep=0pt, label={\color{accent}\small\textbullet}}

% \entry{Title}{Dates}{Organisation, place}
\newcommand{\entry}[3]{%
  \textbf{#1}\hfill{\color{muted}#2}\\
  {\color{muted}\itshape #3}\par\vspace{0.2ex}}

% \pub{Label}{Body}
\newcommand{\pub}[2]{%
  \par\hangindent=0.9cm\hangafter=1\noindent
  \makebox[0.9cm][l]{\color{accent}\bfseries #1}#2\par\vspace{0.8ex}}

\newcommand{\me}{\textbf{Romial Menra}}
\newcommand{\link}[2]{\href{#1}{{\color{accent}#2}}}

% =====================================================================
\begin{document}

% ---- Header ----------------------------------------------------------
\begin{center}
  {\LARGE\bfseries MENRA WEDWANG ROMIAL}\\[0.4ex]
  {\color{muted} Published as \emph{Romial Menra}}\\[0.8ex]
  {\large PhD Candidate in Computer Science, Energy-Aware Cloud Computing}\\[0.3ex]
  {\color{muted} IMT Atlantique, Inria, LS2N (UMR CNRS 6004), Nantes, France}\\[1ex]
  \small
  \faEnvelope\ \link{mailto:itsme@menraromial.com}{itsme@menraromial.com}\quad
  \faGlobe\ \link{https://menraromial.com}{menraromial.com}\quad
  \faGithub\ \link{https://github.com/menraromial}{menraromial}\quad
  \faLinkedin\ \link{https://www.linkedin.com/in/menraromial}{LinkedIn}\\[0.4ex]
  \faOrcid\ \link{https://orcid.org/0009-0007-0943-8593}{0009-0007-0943-8593}\quad
  \faGraduationCap\ \link{https://scholar.google.com/citations?user=M2nDshIAAAAJ}{Google Scholar}\quad
  \faResearchgate\ \link{https://www.researchgate.net/profile/Romial-Menra}{ResearchGate}\quad
  \faLanguage\ French, English
\end{center}

% ---- Profile ---------------------------------------------------------
\section{Research Profile}
PhD candidate working on power management in cloud infrastructures, from hardware-level
power limiting (Intel RAPL, Turbo Boost) to energy-aware orchestration in Kubernetes. My
work combines large-scale experimentation on Grid'5000, analytical modelling of processor
power behaviour, and the design of open-source tools that make power a first-class resource
in the cloud. Before the PhD, two years as a full-stack developer and team lead gave me a
practical, production-oriented view of the systems I now study.

\textbf{Keywords:} power capping, Intel RAPL, energy efficiency, Kubernetes, scheduling,
distributed systems, green IT.

% ---- Education -------------------------------------------------------
\section{Education}
\entry{PhD in Computer Science}{Nov.\ 2024 -- present}{IMT Atlantique, Inria, LS2N, Nantes, France}
Thesis: \emph{Contribution à la gestion de la puissance dans les infrastructures cloud :
de la limitation matérielle à l'orchestration adaptative.}
\vspace{1ex}

\entry{Master of Engineering in Computer Science}{2021 -- 2024}{National Advanced School of Engineering of Yaoundé (ENSPY), Cameroon}
Specialisation in Distributed Systems and Software Engineering.
\vspace{1ex}

\entry{Preparatory Classes, Mathematics \& Physics}{2019 -- 2021}{ENSPY, Cameroon}

% ---- Publications ----------------------------------------------------
\section{Publications}
\textit{International conferences}\par\vspace{0.6ex}
\pub{[C2]}{\me, Guillaume Rosinosky, Remous-Aris Koutsiamanis, Sébastien Bolle, Jean-Marc Menaud.
  \textbf{Understanding Power Limiting Mechanisms in Modern Processors: A Deep Dive Into Intel
  RAPL and Turbo Boost Dynamics.} \emph{Euro-Par 2026: Parallel Processing}, Lecture Notes in
  Computer Science, Springer, pp.\ 135--149, 2026.
  {\color{accent}\bfseries Best Paper Award nominee.} Acceptance rate 26.2\%.\newline
  \link{https://doi.org/10.1007/978-3-032-35251-4_10}{doi:10.1007/978-3-032-35251-4\_10}
  \,·\, \link{https://hal.science/hal-05729597}{hal-05729597}}

\textit{National conferences}\par\vspace{0.6ex}
\pub{[N1]}{\me, Remous-Aris Koutsiamanis, Jean-Marc Menaud.
  \textbf{Intégration de l'Aspect Énergétique dans Kubernetes.} \emph{COMPAS 2024, Conférence
  francophone d'informatique en Parallélisme, Architecture et Système}, Nantes, France, 2024.\newline
  \link{https://hal.science/hal-05638939}{hal-05638939}}

% ---- Talks -----------------------------------------------------------
\section{Talks}
\entry{Understanding Power Limiting Mechanisms in Modern Processors}{Aug.\ 2026}{Euro-Par 2026, 32nd International European Conference on Parallel and Distributed Computing, Pisa, Italy}
\vspace{0.6ex}
\entry{Energy-Aware Kubernetes Orchestration for Grid-Responsive Computing}{Oct.\ 2025}{STACK team seminar, Arzon, France}
\vspace{0.6ex}
\entry{Intégration de l'Aspect Énergétique dans Kubernetes}{Jul.\ 2024}{COMPAS 2024, Nantes, France}

% ---- Experience ------------------------------------------------------
\section{Professional Experience}
\entry{PhD Candidate}{Nov.\ 2024 -- present}{Inria, IMT Atlantique, LS2N, Nantes, France}
\begin{itemize}
  \item Characterised the interaction between Intel RAPL power limits and Turbo Boost on four
        Intel Xeon microarchitectures, and derived an analytical model of turbo duration.
  \item Showed that power capping below about 50\% of TDP systematically degrades energy efficiency.
  \item Designing Kubernetes components for power-aware operation: node-level power capping,
        automatic server power profiling and energy-aware scheduling.
  \item Large-scale experimentation on the Grid'5000 testbed; collaboration with Orange Research.
\end{itemize}

\entry{Research Intern}{Apr.\ 2024 -- Aug.\ 2024}{IMT Atlantique, Nantes, France}
\begin{itemize}
  \item Implemented server power limitation with Intel RAPL and evaluated it on Grid'5000.
  \item Developed an energy-aware Kubernetes scheduler with custom resource constraints.
  \item Results published at COMPAS 2024.
\end{itemize}

\entry{Software Developer \& Team Lead}{May 2022 -- Mar.\ 2024}{Forall Founders, Cameroon}
\begin{itemize}
  \item Led a team of 10 developers building e-learning platforms and school management systems
        (Django, Node.js, React, Next.js, Vue.js).
  \item Project manager of a mobile e-learning application with 15,000+ active users
        (\link{https://tefaafrik.com}{TEFAA}); cut latency from 5\,s to 1.2\,s through query
        optimisation and caching.
  \item Built an IoT attendance system (PyQt, Raspberry Pi, Arduino) and CI/CD pipelines on Linode.
\end{itemize}

% ---- Teaching & supervision -----------------------------------------
\section{Teaching \& Supervision}
\entry{Teaching Assistant, Databases (TD/TP)}{Spring 2025}{IMT Atlantique, Nantes, first-year engineering students}
Relational model, SQL, normalisation, UML conceptual modelling, programmatic interfaces.
\vspace{1ex}

\entry{Teaching Assistant, Cloud Computing with VMware vSphere (TP)}{Spring 2025}{IMT Atlantique, Nantes}
\vspace{0.6ex}

\entry{Co-supervision, Clement Obama (M2 research intern)}{Apr.\ 2026 -- Sep.\ 2026}{IMT Atlantique, Nantes}
\emph{Reactive control of energy constraints in Kubernetes clusters: multi-lever
re-adaptation policy with SLA management.}

% ---- Service & community ---------------------------------------------
\section{Service \& Community}
\begin{itemize}
  \item Student volunteer, UCC \& BDCAT 2025 (Utility and Cloud Computing / Big Data
        Computing, Applications and Technologies), Nantes, Dec.\ 2025.
  \item Founder of \textbf{C5IN}, Cameroon Cloud-Edge-IoT Innovation Network.
  \item Contributor to \textbf{DIBAC}, an open-access institutional repository for Cameroonian
        research, and \textbf{TupuriHub}, a digital hub for Toupouri culture and heritage.
  \item Technical blog on Kubernetes and Linux power management at \link{https://menraromial.com}{menraromial.com}.
\end{itemize}

% ---- Skills ----------------------------------------------------------
\section{Technical Skills}
\begin{tabularx}{\linewidth}{@{}>{\bfseries}p{3.4cm}X@{}}
  Research \& HPC     & Grid'5000, Intel RAPL, powercap, energy profiling, performance benchmarking \\
  Cloud \& DevOps     & Kubernetes (k8s, k3s, kubeadm, Kubebuilder, scheduler plugins), Docker,
                        Terraform, Ansible, GitHub Actions, GitLab CI, Jenkins, Prometheus, Grafana \\
  Languages           & Go, Python, C, Java, JavaScript/TypeScript, Bash \\
  Web                 & Django, FastAPI, Node.js, React, Next.js, Vue.js, Spring Boot \\
  Data                & PostgreSQL, MySQL, MongoDB, Redis, Pandas, NumPy \\
\end{tabularx}

% ---- Certifications & interests --------------------------------------
\section{Certifications \& Interests}
\textbf{Certifications:} IBM (Coursera, 2023): Agile Development and Scrum, Applied Data Science
Capstone, Machine Learning with Python, Data Visualization with Python, Databases and SQL for
Data Science. EPFL (Coursera, 2021): Programming and Object-Oriented Programming in Java.

\vspace{0.6ex}
\textbf{Interests:} open source, technical writing, reading (science, leadership, history),
documentaries, cooking, cinema, nature.

\vfill
{\centering\footnotesize\color{muted} Last updated: September 2026\par}

\end{document}
